\section{\npcTaravangian}
\label{sec:taravangian}

In this chapter, the party will meet \npcTaravangian{} twice - once when they petition for an access to the \locationPalanaeum{} and the second one at the end of their descent into the Conclave Hospital.

\reasonBox{
	When the party enters the Conversation with \enemyTaravangian{} for the second time, make sure his dominant attribute is different than before so the players could experience both sides of the character. In any subsequent Conversations you can allow any configuration of his attributes.
}

\reasonBox{
	Due to the cannon lore, please try to make sure the king survives this chapter. In case of any danger, allow \enemyMalata{} to join the fight and defend him as he flees the scene.
}




\subsection{Petition to the King}
\label{ssec:taravangianPetition}

When the party attempts to receive access to enter the \locationPalanaeum{}, read the following:

\quoteBox{
	You enter the only palace wing opened to the public. Today, \npcTaravangian{} listens to the pleas of his subjects from a high stone balcony. \\ \\
	When it is your turn to present your case, the kings welcomes you: "Ah, our honored guests from Alethkar. What boon do you require?"
}

To obtain access to the \locationPalanaeum{}, the party will need to enter the Conversation with \enemyTaravangian{}. Roll for his \abilityTaravangianBoonName{} ability to determine his attributes, Cognitive and Spiritual Defenses and total ranks of his skills.

The king aims for the greater good of all humanity and he sees the PCs as useful tools he could recruit later. However, during a public event as this, he will refrain from any open moves and he will allow the party a bit more ignorance.

It is not predetermined what arguments will resonate with him the most. Due to his \abilityTaravangianBoonName{} ability, his personality changes and so does his behavior in the Conversation.

\conversationEnd{\enemyTaravangian{}}{1}

\subsubsection{King's Guards}

\enemyTaravangian{} is surrounded by four \enemyDiagramMuscle{}s. Behind him, hidden in shadows, \enemyMalata{} watches his back. They do not participate in the Conversation in any way, but they will defend the king in case of any violence from the party.

\subsubsection{Level Up!}

Regardless of the result of the Conversation with the king, advance the party to level 9 once they complete this scene.

\subsubsection{Clues}

If the party convinces \npcTaravangian{} to give them access to the library, they will be able to enter inside  \fullref{sec:palanaeum} and from there, with some scouting - even to \fullref{sec:conclaveHospital}.

If the party fails, they will need to find some other way of entering \fullref{sec:conclaveHospital}. Potential options may include \fullref{ssec:cityHospitalConclaveTransfer}, or following the surgeon or the soulcaster to find their way in, but you can allow for any other reasonable and creative solutions.



\newpage
\subsection{Leader of the Diagram}
\label{ssec:taravangianDiagram}

When you start this scene, read the following:

\quoteBox{
	\npcTaravangian{} is standing on the balcony overseeing the Conclave Hospital. There are guards all around him. \\ \\
	"And the vision? It must have been about \locationUrithiru{}. We can ship to Vedenar to get there.", he says.
}

The party is about to enter a Conversation with \enemyTaravangian{}. This, however, may look differently based on the party's choice in \fullref{ssec:conclaveHospitalPatient} scene.

\subsubsection{Unfortunate Accident}

If the party talked \npcDiagramSurgeon{} into killing himself, the king will assume their hostility at first. If they wish to join the Diagram, they would need to convince him about their true intentions and honest regret about the surgeon.

\subsubsection{Opposing the Diagram}

If the party killed \npcDiagramSurgeon{}, \npcTaravangian{} will see them as a potential threat to everything he is working on. The party must convince him not to send his forces against them and let them leave with \npcGhostbloodWounded{}.

\subsubsection{Joining the Diagram}

If the party killed \npcGhostbloodWounded{}, \npcTaravangian{} will gladly welcome them to his organization and allow to travel with him to \locationUrithiru{}. However, there might still be some terms and conditions of this arrangement that the party might wish to negotiate.

\subsubsection{The Conversation}

Before the Conversation, determine attributes of the king that are affected by his \abilityTaravangianBoonName{}. The king's goals and motivation may differ depending on the PCs' choices and on his current level of intelligence.

\conversationEnd{\enemyTaravangian{}}{1}

\reasonBox{
	The whole point of having two Conversation scenes with \npcTaravangian{} is to show his dual aspect and have two different experiences with the same NPC. Make sure the second scene with him is using different dominant attribute than the first one, switch the attributes if required.
}

\subsubsection{Aftermath}

If the party succeeded on the Conversation, their request will be accepted and honored by \npcTaravangian{}.

If they fail, their request will be denied. If that request was to allow them leave, the king will send six \enemyDiagramMuscle{}s and \enemyMalata{} to slay the party. Also, the \enemyDiagramScribe{} from the previous scene will join to support her allies.

\reasonBox{
	If the guards are sent to kill the party, allow \enemyMalata{} to show off her skills at the moments before the battle (instantly summoning the Shardblade, using Division on the nearby stones for dramatic effect, etc). That will allow the players to recognize the threat and flee if they believe it too risky. In that case, don't pursue them.
} 